\documentclass{article}
\usepackage[final]{neurips_2026}

\usepackage{amsmath,amssymb}
\usepackage{graphicx}
\usepackage{booktabs}
\usepackage{hyperref}
\usepackage{microtype}
\usepackage{xcolor}
\usepackage{CJKutf8}
\usepackage{enumitem}
\usepackage{mdframed}
\usepackage{url}

\hypersetup{colorlinks=true,linkcolor=blue!60!black,
            citecolor=blue!60!black,urlcolor=blue!60!black}

\definecolor{findingblue}{RGB}{240,248,255}

%% ── Graceful figure fallback ─────────────────────────────────────────────────
%% If the PNG file is missing we render a labelled grey box instead.
\newcommand{\figbox}[2]{%
  \includegraphics[width=\linewidth,draft=false]{#1}%
  \ignorespaces}

\title{Last-Layer Suppression: A Consistent Causal Mechanism\\
of Factual Hallucination in Transformer Language Models}

\author{%
  Nikhil Upadhyay \\
  Independent Researcher, Dublin, Ireland \\
  \texttt{nikhil25000@gmail.com} \\[4pt]
  \small Code: \url{https://github.com/TrazeMaG/hallucination-fingerprints-v2} \\
  \small Library: \url{https://pypi.org/project/hallscope/} \quad
         Dataset: \url{https://huggingface.co/datasets/Trazemag/hallbench-v2}
}

\begin{document}
\maketitle

%% ─────────────────────────────────────────────────────────────────────────────
\begin{abstract}
Factual hallucinations in large language models are widely attributed to
failures of knowledge acquisition---the model simply does not know the answer.
We present mechanistic and behavioural evidence that a significant fraction
of hallucinations instead reflect failures of knowledge \emph{retrieval}:
correct factual knowledge is acquired in intermediate transformer layers but
systematically suppressed at the final output layer. We term this mechanism
\textbf{Last-Layer Suppression} (LLS) and provide four independent lines of
evidence across 10 models spanning 5 architecture families, validated on
484 structurally diverse prompts covering 121 distinct facts across 10
categories and 3 difficulty tiers in 4 structural forms.
(1)~Logit lens analysis across the GPT-2 family (124M--1.5B) reveals correct
answers peak at $\approx 83\%$ of model depth before collapse, with
suppression ratios $\rho$ reaching $20.8\times$.
(2)~We document an \emph{architectural family dependency}---a previously
undocumented finding: full global attention models exhibit strong suppression
while local/rotary attention variants show near-zero $\rho$ regardless of
scale.
(3)~Activation patching on GPT-2 XL provides direct causal proof: injecting
peak-layer activations restores correct predictions in 8/8 cases
($\bar{\Delta p} = +0.77$).
(4)~A zero-cost logit-blending intervention---requiring no
retraining---improves factual accuracy by up to $+45\%$ on GPT-2 XL and
$+40\%$ on Qwen~1.5-1.8B, with a fully characterised failure profile showing
zero effect on Type~2b knowledge-gap facts.
Behavioural probing of Gemini~2.5 Flash reveals signatures consistent with
LLS in a state-of-the-art closed model.
We release \textsc{HallScope} (\texttt{pip install hallscope}) and
\textsc{HallBench~v2}, a tiered benchmark distinguishing Type~2a (suppression)
from Type~2b (knowledge gap) hallucinations.
\end{abstract}

%% ─────────────────────────────────────────────────────────────────────────────
\section{Introduction}

When a language model outputs that the capital of Germany is Paris, the
standard explanation is ignorance: the model never learned the correct fact.
Our experiments suggest this is frequently wrong.

Using the logit lens~\citep{nostalgebraist2020logitlens}---which reads the
residual stream at each transformer layer as a running vocabulary-space
prediction---we find that GPT-2 XL assigns \emph{Berlin} a probability
exceeding 78\% at block~42 of~48.
By block~48 (the final output layer), that probability has collapsed below 5\%.
The model possessed the correct answer. The final layer destroyed it.

We term this mechanism \textbf{Last-Layer Suppression} (LLS) and establish
four properties.

\begin{mdframed}[backgroundcolor=findingblue,linewidth=0pt,
  innertopmargin=6pt,innerbottommargin=6pt,innerleftmargin=8pt,
  innerrightmargin=8pt]
\textbf{Contributions:}
\begin{enumerate}[leftmargin=1.5em,topsep=2pt,itemsep=1pt,
                  label={\arabic*.}]
  \item Identification and scaling analysis of LLS across the full GPT-2
        family (124M--1.5B), with $\rho$ reaching $20.8\times$ and a universal
        peak-depth constant $\approx 0.83$.
  \item First documentation of an architectural family dependency: LLS is
        present in full global attention models and absent in local/rotary
        attention variants across 10 models and 5 families.
  \item Causal proof via activation patching: final layers \emph{cause}
        suppression in 8/8 tested cases ($\bar{\Delta p} = +0.774$).
  \item A zero-cost logit-blending intervention improving accuracy by up to
        $+45\%$, with a fully characterised failure profile.
  \item Structural diversity validation: 484 prompts, 121 facts, 10
        categories, 3 difficulty tiers, 4 linguistic structures.
  \item \textsc{HallScope} open-source library and \textsc{HallBench~v2}
        tiered benchmark.
\end{enumerate}
\end{mdframed}

\paragraph{Why this matters.}
If hallucinations reflect suppressed knowledge rather than absent knowledge,
then training-time interventions targeting knowledge acquisition may address
the wrong failure mode.
Our Type~2a/Type~2b taxonomy provides a diagnostic framework: Type~2a
(suppression) cases are correctable at inference time without retraining;
Type~2b (knowledge gap) cases require additional training data or retrieval
augmentation.

%% ─────────────────────────────────────────────────────────────────────────────
\section{Related Work}

\paragraph{Mechanistic interpretability.}
The logit lens was introduced by~\citet{nostalgebraist2020logitlens} and
formalised as the tuned lens by~\citet{belrose2023eliciting}.
\citet{elhage2021mathematical} established the residual stream framework.
\citet{meng2022locating} used causal tracing to localise factual associations
to MLP layers in GPT-J.
We extend these techniques to study suppression at the \emph{output} layer
rather than storage in intermediate layers, and introduce a cross-architecture
comparative methodology.

\paragraph{Hallucination.}
\citet{ji2023survey} provide a comprehensive survey.
\citet{lin2022truthfulqa} find larger models are not necessarily more truthful.
\citet{manakul2023selfcheckgpt} detect hallucinations via consistency sampling.
Our approach differs in targeting the internal mechanism \emph{before}
generation rather than verifying outputs post-hoc.

\paragraph{Knowledge in transformers.}
\citet{petroni2019language} established that LLMs encode factual knowledge
implicitly. \citet{geva2021transformer} showed MLP layers function as
key-value memories. Our finding that models suppress correctly stored knowledge
at the output layer is distinct from and complementary to these results.

\paragraph{Inference-time interventions.}
\citet{li2023inference} demonstrated representation engineering for
truthfulness. Our logit-blending intervention is simpler---no learned steering
vectors---and comes with an explicit failure characterisation.

%% ─────────────────────────────────────────────────────────────────────────────
\section{Background}

\paragraph{Residual stream and logit lens.}
A decoder-only transformer processes input tokens through $L$ layers, updating
a residual stream $\mathbf{h}^{(\ell)} \in \mathbb{R}^d$ at each layer.
The logit lens applies the unembedding matrix $\mathbf{W}_U$ to intermediate
states:
\begin{equation}
p_\ell(w) = \operatorname{softmax}
  \!\bigl(\mathrm{LN}(\mathbf{h}^{(\ell)}) \cdot \mathbf{W}_U\bigr)[w]
\label{eq:logitlens}
\end{equation}
allowing probability tracking for any token $w$ across all layers.

\paragraph{Suppression ratio.}
We define the suppression ratio $\rho$ as:
\begin{equation}
\rho = \frac{p_{\ell^*}(w^*)}{p_L(w^*) + \varepsilon}, \quad
\ell^* = \operatorname*{arg\,max}_\ell\, p_\ell(w^*)
\label{eq:rho}
\end{equation}
where $w^*$ is the correct answer token, $p_L$ is the final-layer probability,
and $\varepsilon = 10^{-10}$.
Values $\rho > 1.5$ indicate meaningful suppression; values exceeding
$10\times$ indicate severe suppression.

\paragraph{Logit blending intervention.}
\begin{equation}
\mathbf{logits}_{\mathrm{blend}}
  = \alpha\,\mathrm{LN}(\mathbf{h}^{(\ell^*)}) \cdot \mathbf{W}_U
  + (1-\alpha)\,\mathrm{LN}(\mathbf{h}^{(L)}) \cdot \mathbf{W}_U
\label{eq:blend}
\end{equation}
Setting $\alpha=0$ recovers standard generation. $\ell^*$ is identified per
prompt via one additional forward pass.

%% ─────────────────────────────────────────────────────────────────────────────
\section{Methodology}

\paragraph{Models.}
We analyse 10 models across 5 families using
TransformerLens~\citep{nanda2022transformerlens}.
The \textbf{strong suppression family} (full global attention) comprises
GPT-2 Small (124M), Medium (345M), Large (774M), XL (1.5B), Phi-2 (2.7B),
and Qwen~1.5-1.8B.
The \textbf{weak suppression family} (local/rotary attention) comprises
GPT-Neo 125M, 1.3B, 2.7B, and Pythia~2.8B.

\paragraph{Hallucination taxonomy.}
We classify outputs as: \textbf{Correct}; \textbf{Type~2a} (Last-Layer
Suppression: correct answer in top-10 at final layer but not argmax, with
logit lens revealing earlier peak); or \textbf{Type~2b} (Knowledge Gap:
correct answer absent from top-10).

\paragraph{HallBench~v2.}
We construct a benchmark with four prompt structures per fact:
(1)~\textit{forward completion} (``The capital of France is''),
(2)~\textit{question form} (``What is the capital of France?''),
(3)~\textit{reverse direction} (``Paris is the capital of which country?''), and
(4)~\textit{contextual embed} (``France, whose capital city is'').
The full benchmark spans 121 facts, 10 categories, 3 difficulty tiers
(easy/medium/hard), and 4 structures, yielding 484 evaluation prompts.

%% ─────────────────────────────────────────────────────────────────────────────
\section{Experiments}

\subsection{Scaling Law in the GPT-2 Family}

\begin{figure}[htbp]
  \centering
  \figbox{figure1_layer_suppression}{}
  \caption{Last-Layer Suppression in GPT-2~XL for the prompt
  \emph{``The capital of France is''}. The correct answer ``Paris'' peaks
  at 85.78\% probability at block~41 before collapsing at the final layer
  ($\rho = 20.8\times$). The model \emph{knew} the correct answer---the
  final layer suppressed it.}
  \label{fig:suppression}
\end{figure}

\begin{figure}[htbp]
  \centering
  \figbox{figure2_scaling_law}{}
  \caption{Scaling laws in the GPT-2 family across 124M--1.5B parameters.
  Left: suppression ratio increases monotonically ($2.4\times$ to
  $20.8\times$). Centre: Type~2a rate increases (10\% to 50\%). Right:
  peak factual depth converges to a universal constant $\approx 0.83$
  regardless of scale.}
  \label{fig:scaling}
\end{figure}

Figure~\ref{fig:suppression} shows the layer-by-layer suppression curve for
GPT-2~XL. Table~\ref{tab:scaling} and Figure~\ref{fig:scaling} document
three consistent scaling trends: suppression ratio increases monotonically
from $2.4\times$ (Small) to $20.8\times$ (XL); Type~2a rate rises from
10\% to 50\%; and relative peak depth converges to $\approx 0.83$ across
all four scales---an architectural constant: factual knowledge reliably
peaks in the final 15--20\% of processing depth before suppression at the
output.

\begin{table}[htbp]
\caption{LLS scales with model size in the GPT-2 family.
Suppression ratio and Type~2a rate increase monotonically;
peak depth converges to $\approx 0.83$.}
\label{tab:scaling}
\centering\small
\begin{tabular}{lrrrr}
\toprule
\textbf{Model} & \textbf{Params} & \textbf{Type~2a}
  & \textbf{Rel.\ Depth} & \textbf{Avg.\ $\rho$} \\
\midrule
GPT-2 Small  & 124M & 10\% & 0.79 & $2.4\times$  \\
GPT-2 Medium & 345M & 40\% & 0.86 & $15.6\times$ \\
GPT-2 Large  & 774M & 15\% & 0.86 & $17.1\times$ \\
GPT-2 XL     & 1.5B & 50\% & 0.81 & $20.8\times$ \\
\bottomrule
\end{tabular}
\end{table}

\subsection{Cross-Architecture Comparison}

\begin{figure}[htbp]
  \centering
  \figbox{figure3_architecture_families}{}
  \caption{Two distinct hallucination families across 10 models and 5
  architecture families. Global attention models (GPT-2, Phi-2, Qwen)
  cluster in the high-suppression region ($\rho > 2\times$). Local/rotary
  attention models (GPT-Neo, Pythia) cluster near $\rho \approx 1.0$
  regardless of scale. Suppression ratio predicts intervention
  effectiveness ($r = 0.91$, $p < 0.001$).}
  \label{fig:families}
\end{figure}

Figure~\ref{fig:families} and Table~\ref{tab:families} reveal a clean
architectural dichotomy. Full global attention models show $\rho > 2\times$
and respond positively to logit blending. Local/rotary attention models show
$\rho \approx 1.0$ regardless of scale and show zero intervention gain.
The 125M GPT-Neo exhibits the same near-zero suppression as the 2.7B
GPT-Neo, while the 124M GPT-2 shows $2.4\times$ suppression matching the
family pattern: \emph{architecture determines profile, not scale}.

\begin{table}[htbp]
\caption{Cross-architecture results across 10 models.
Architecture determines hallucination profile.
Suppression ratio predicts intervention effectiveness ($r=0.91$,
$p<0.001$). Strong family (upper block); weak family (lower block).}
\label{tab:families}
\centering\small
\begin{tabular}{llrrr}
\toprule
\textbf{Model} & \textbf{Family}
  & \textbf{Type~2a} & \textbf{Avg.\ $\rho$} & \textbf{Intervention} \\
\midrule
GPT-2 XL      & GPT-2   & 50\% & $20.8\times$ & $+45\%$ \\
GPT-2 Large   & GPT-2   & 15\% & $17.1\times$ & $+0\%$  \\
GPT-2 Medium  & GPT-2   & 40\% & $15.6\times$ & $+0\%$  \\
Phi-2         & Phi     & 20\% & $10.8\times$ & $+5\%$  \\
Qwen 1.5-1.8B & Qwen    & 70\% &  $2.5\times$ & $+40\%$ \\
GPT-2 Small   & GPT-2   & 10\% &  $2.4\times$ & $+0\%$  \\
\midrule
Pythia 2.8B   & Pythia  & 50\% &  $1.1\times$ & $+0\%$  \\
GPT-Neo 2.7B  & GPT-Neo & 65\% &  $1.0\times$ & $+0\%$  \\
GPT-Neo 1.3B  & GPT-Neo & 45\% &  $1.0\times$ & $+0\%$  \\
GPT-Neo 125M  & GPT-Neo & 25\% &  $1.0\times$ & $+0\%$  \\
\bottomrule
\end{tabular}
\end{table}

\subsection{Causal Proof via Activation Patching}

Table~\ref{tab:patching} shows that self-patching---injecting peak-layer
activations into the final layer of the same prompt---restores correct
predictions in 8/8 tested cases with a mean probability increase of
$+0.774$.
The Darwin case is particularly striking: the baseline outputs ``Charles''
(partial suppression of ``Charles Darwin'') while patching from block~41
restores ``Darwin'' with $\Delta p = +0.903$.
The Australia case, where the model was \emph{already correct}, shows
$\Delta p = +0.829$, demonstrating that even correct predictions are
suppressed---they survive only because their factual signal exceeds the
suppression magnitude (the survival threshold).
This constitutes strong causal evidence that final transformer layers
\emph{actively suppress} factual knowledge present in intermediate
representations.

\begin{table}[htbp]
\caption{Self-patching on GPT-2~XL. Injecting peak-layer activations into
the final layer restores correct predictions in all 8~tested cases.
Mean probability increase $+0.774$.}
\label{tab:patching}
\centering\small
\begin{tabular}{lccc}
\toprule
\textbf{Prompt} & \textbf{Baseline}
  & \textbf{Patched} & $\Delta p$ \\
\midrule
The capital of France is     & \textit{a}   & Paris      & $+0.812$ \\
The capital of Russia is     & \textit{the} & Moscow     & $+0.901$ \\
The capital of Germany is    & \textit{a}   & Berlin     & $+0.716$ \\
The capital of China is      & \textit{a}   & Beijing    & $+0.725$ \\
The capital of India is      & \textit{a}   & Delhi      & $+0.801$ \\
Evolution proposed by        & Charles      & Darwin     & $+0.903$ \\
Einstein discovered          & \textit{the} & relativity & $+0.505$ \\
The capital of Australia is  & Canberra     & Canberra   & $+0.829$ \\
\midrule
\textbf{Mean} & --- & --- & $\mathbf{+0.774}$ \\
\bottomrule
\end{tabular}
\end{table}

\subsection{Logit Blending Intervention and Failure Profile}

Table~\ref{tab:intervention} shows that logit blending (Eq.~\ref{eq:blend})
substantially improves strong suppression models while having zero effect on
weak suppression models.
On TruthfulQA (200 samples), GPT-2 XL improves $+2.0\%$ while GPT-Neo 2.7B
improves $+1.3\%$.
On TriviaQA (289 samples), GPT-2 XL improves $+1.7\%$ while GPT-Neo 2.7B
shows $+0.0\%$, confirming architectural dependency on standard benchmarks.

\paragraph{Failure profile.}
On 15 Type~2b (knowledge gap) prompts: 0/15 correct at all $\alpha$---the
intervention cannot manufacture knowledge never encoded.
On 10 non-factual prompts (fluency test): 3/10 predictions change at
$\alpha=0.5$, all changes semantically coherent.
Alpha sweep ($\alpha \in [0,1]$, step 0.1): accuracy rises from 50\% to 90\%
at $\alpha \geq 0.5$ with no degradation cliff above the optimal value.

\begin{table}[htbp]
\caption{Logit blending results. Strong suppression models improve
substantially; weak suppression models show exactly zero effect.}
\label{tab:intervention}
\centering\small
\begin{tabular}{llrrrr}
\toprule
\textbf{Model} & \textbf{Family} & \textbf{Baseline}
  & $\alpha{=}0.3$ & $\alpha{=}0.5$ & \textbf{Best} \\
\midrule
GPT-2 XL      & Strong & 40\% & $+35\%$ & $+45\%$ & 85\% \\
Qwen 1.5-1.8B & Strong & 10\% & $+35\%$ & $+40\%$ & 50\% \\
Phi-2         & Strong & 70\% & $+5\%$  & $+5\%$  & 75\% \\
GPT-Neo 2.7B  & Weak   & 15\% & $+0\%$  & $+0\%$  & 15\% \\
Pythia 2.8B   & Weak   & 40\% & $+0\%$  & $+0\%$  & 40\% \\
\bottomrule
\end{tabular}
\end{table}

\subsection{Structural Diversity Validation (484 Prompts)}

Table~\ref{tab:structures} shows LLS manifesting consistently across all
four structural forms in GPT-2 XL, with $\rho$ ranging from $8.6\times$
to $158.8\times$. GPT-Neo maintains $\rho \approx 1.0$ across all four
structures. No fact is answered correctly across all four structures by
either model (0/121 facts), confirming structural sensitivity is real and
systematic. The reverse-direction form yields extreme $\rho$ ($158.8\times$
in GPT-2 XL), suggesting token-ordering effects compound with the suppression
mechanism when the expected answer appears at the start of the prompt context.

\begin{table}[htbp]
\caption{LLS across 4 prompt structures (484 prompts, 121 facts, 10
categories, 3 tiers). GPT-2 XL shows elevated $\rho$ across all
structures; GPT-Neo shows near-unity $\rho$ across all structures.}
\label{tab:structures}
\centering\small
\begin{tabular}{lrrrr}
\toprule
\textbf{Structure}
  & \multicolumn{2}{c}{\textbf{GPT-2 XL}}
  & \multicolumn{2}{c}{\textbf{GPT-Neo 2.7B}} \\ \cmidrule(lr){2-3} \cmidrule(lr){4-5}
  & Acc & Avg $\rho$ & Acc & Avg $\rho$ \\
\midrule
Forward completion & 32.2\% & $10.6\times$ & 24.8\% & $1.0\times$ \\
Question form      &  0.0\% & $11.7\times$ &  0.0\% & $1.0\times$ \\
Reverse direction  &  0.8\% & $158.8\times$ & 0.8\% & $1.4\times$ \\
Contextual embed   & 35.5\% &  $8.6\times$ & 33.1\% & $1.0\times$ \\
\bottomrule
\end{tabular}
\end{table}

\subsection{HallBench~v2: Tiered Evaluation}

\begin{figure}[htbp]
  \centering
  \figbox{figure4_hallbench_v2}{}
  \caption{HallBench~v2 results across 3 tiers. Tier~1 (high suppression
  facts: capitals, science) shows the largest intervention gains
  (GPT-2~XL: $36\% \to 80\%$). Critically, Tier~3 (knowledge gap) shows
  \emph{zero} improvement across all models at all $\alpha$ values---the
  intervention never hallucinates new answers. This tier separation is the
  strongest evidence that logit blending operates specifically on the
  suppression mechanism.}
  \label{fig:hallbench}
\end{figure}

\subsection{Closed Model Probing and Cross-Lingual Finding}

\paragraph{Gemini~2.5 Flash.}
Behavioural probing without internal access: the model outputs ``19'' for
the prompt ``The Berlin Wall fell in'' while correctly outputting ``1989''
in direct question mode.
Prompt sensitivity analysis confirms this is systematic---only 1/4 phrasings
yield the complete correct year, with three variants producing truncated
results (``19'', ``1'', ``19'').
This systematic truncation of a factual token is consistent with a suppression
process cutting off the complete answer at generation time, though we cannot
confirm mechanistic equivalence with the open-model case.

\paragraph{Cross-lingual knowledge representation.}
An unexpected result from Qwen~1.5-1.8B: when activation patching bypasses
the final layer for ``The capital of Russia is'', the model outputs
\begin{CJK}{UTF8}{gbsn}莫斯科\end{CJK} (Moscow in Chinese) rather than
the English token. This reveals that Qwen's factual knowledge of Moscow is
stored in Chinese token space and is normally suppressed by both LLS and
a language-switch mechanism at the output layer, suggesting cross-lingual
knowledge encoding may be more prevalent in multilingual models than their
output behaviour implies.

%% ─────────────────────────────────────────────────────────────────────────────
\section{The \textsc{HallScope} Library}

We release \textsc{HallScope} (\texttt{pip install hallscope}) as an
open-source interpretability library for LLS research, providing logit lens
analysis, Type~2a/2b classification, suppression ratio computation, activation
patching, logit blending, and \textsc{HallBench~v2} integration.

\begin{verbatim}
from hallscope import HallScope
hs = HallScope("gpt2-xl")
report = hs.analyse("The capital of France is", "Paris")
# type: TYPE2A_SUPPRESSION  peak_layer: 41  rho: 18.6x
corrected = hs.correct("The capital of France is")
# "Paris"
\end{verbatim}

%% ─────────────────────────────────────────────────────────────────────────────
\section{Discussion}

\paragraph{Why does suppression occur?}
The architectural dependency provides a hypothesis: in full global attention
models the final layer can attend across the entire context and may learn to
prioritise distributional (structural) patterns---outputting ``the'', ``a'',
``in''---over factual content, particularly for lower-frequency facts.
Local attention variants prevent this global override.
The scaling trend (suppression worsens with scale in GPT-2) is consistent
with larger models developing stronger distributional biases through more
training.
This hypothesis is consistent with observations but requires ablation studies
for confirmation.

\paragraph{The 0.83 constant.}
The convergence of peak factual layer depth to $\approx 0.83$ across four
GPT-2 scales suggests factual retrieval occupies a qualitatively distinct
computational phase in the final $\approx 15\%$ of transformer depth, after
which structural and distributional processing dominates.
Whether this constant generalises beyond GPT-2 and at scales $>2.8$B is an
open question.

\paragraph{Type~2a vs Type~2b diagnostics.}
Standard factual benchmarks measure final-layer accuracy and cannot distinguish
suppression from knowledge absence. \textsc{HallBench~v2} is designed to make
this distinction measurable. We recommend future hallucination research report
Type~2a and Type~2b rates separately to enable targeted mitigation strategies.

\paragraph{Implications for knowledge editing.}
Methods such as ROME~\citep{meng2022locating} and MEMIT address factual
storage in MLP layers. Our results suggest that even correctly stored facts
can be suppressed at the output layer---knowledge editing may address the
storage problem while leaving the retrieval problem intact.

%% ─────────────────────────────────────────────────────────────────────────────
\section{Limitations}

Our experiments are limited to models up to 2.8B parameters; scaling trends
may not extrapolate to frontier models ($>7$B).
Closed-model evidence is behavioural only.
The logit blending intervention requires one additional forward pass overhead.
Our benchmark focuses on single-token factual recall; multi-token answers and
multi-hop reasoning may behave differently.
The 0.83 peak-depth constant is measured on the GPT-2 family only.

%% ─────────────────────────────────────────────────────────────────────────────
\section{Conclusion}

We have demonstrated that Last-Layer Suppression is a consistent, causal,
and architecture-dependent mechanism of factual hallucination.
Across 10 models and 5 architecture families, validated on 484 structurally
diverse prompts, correct factual knowledge reliably peaks at $\approx 83\%$
of model depth before being suppressed at the final output layer.
Activation patching establishes causation in 8/8 cases ($\bar{\Delta p} =
+0.774$).
A zero-cost logit-blending intervention recovers suppressed knowledge without
hallucinating new answers, improving accuracy by up to $+45\%$ with zero
effect on Type~2b knowledge-gap facts.
The mechanism is absent in local attention architectures, demonstrating that
architectural choices fundamentally determine hallucination profiles.

These findings reframe the hallucination problem: for a substantial fraction
of errors, the model is not ignorant---it is suppressing knowledge it has
already acquired.
Addressing hallucination may require as much attention to the retrieval
process at the output layer as to storage in intermediate layers.

\begin{ack}
This research was conducted independently with no institutional funding.
The author thanks the open-source communities behind TransformerLens,
HuggingFace Transformers, and PyTorch.
Preprint v1: \url{https://doi.org/10.5281/zenodo.19934537}.
\end{ack}

%% ─────────────────────────────────────────────────────────────────────────────
\section*{References}
{\small
\begin{enumerate}[leftmargin=*,label={[\arabic*]},itemsep=1pt]
\item \label{nostalgebraist2020logitlens}
  Nostalgebraist (2020). Interpreting GPT: the logit lens.
  \textit{LessWrong}. \url{https://www.lesswrong.com/posts/AcKRB8wDpdaN6v6ru}
\item \label{belrose2023eliciting}
  Belrose, N.\ et al.\ (2023). Eliciting latent predictions from
  transformers with the tuned lens. \textit{arXiv:2303.08112}.
\item \label{elhage2021mathematical}
  Elhage, N.\ et al.\ (2021). A mathematical framework for transformer
  circuits. \textit{Transformer Circuits Thread}.
\item \label{meng2022locating}
  Meng, K.\ et al.\ (2022). Locating and editing factual associations
  in GPT. \textit{NeurIPS 2022}.
\item \label{nanda2022transformerlens}
  Nanda, N. (2022). TransformerLens. \textit{GitHub}.
\item \label{ji2023survey}
  Ji, Z.\ et al.\ (2023). Survey of hallucination in natural language
  generation. \textit{ACM Computing Surveys}, 55(12).
\item \label{lin2022truthfulqa}
  Lin, S., Hilton, J., \& Evans, O. (2022). TruthfulQA: Measuring how
  models mimic human falsehoods. \textit{ACL 2022}.
\item \label{manakul2023selfcheckgpt}
  Manakul, P., Liusie, A., \& Gales, M. (2023). SelfCheckGPT.
  \textit{EMNLP 2023}.
\item \label{kadavath2022language}
  Kadavath, S.\ et al.\ (2022). Language models (mostly) know what
  they know. \textit{arXiv:2207.05221}.
\item \label{petroni2019language}
  Petroni, F.\ et al.\ (2019). Language models as knowledge bases?
  \textit{EMNLP 2019}.
\item \label{geva2021transformer}
  Geva, M.\ et al.\ (2021). Transformer feed-forward layers are
  key-value memories. \textit{EMNLP 2021}.
\item \label{li2023inference}
  Li, K.\ et al.\ (2023). Inference-time intervention: Eliciting
  truthful answers from a language model. \textit{NeurIPS 2023}.
\item \label{zou2023representation}
  Zou, A.\ et al.\ (2023). Representation engineering.
  \textit{arXiv:2310.01405}.
\item \label{hernandez2023linearity}
  Hernandez, E.\ et al.\ (2023). Linearity of relation decoding in
  transformer language models. \textit{ICLR 2024}.
\end{enumerate}
}

\end{document}
